\newif\ifglobal

%%%%%%%%%%%%%%%%%%%%%%%
%%% Compile to view %%%
%%%%%%%%%%%%%%%%%%%%%%%

%%%%%%%%%%%%%%%%%%%%%%%%%%%%%%%%%%%%%%%%%%%%%%%%%%%
%%% Readme：                                    %%%
%%% https://www.overleaf.com/read/bwdjvfjksvjy/ %%%
%%%%%%%%%%%%%%%%%%%%%%%%%%%%%%%%%%%%%%%%%%%%%%%%%%%

% >>> M?: marks a module setting number or important notice

% >>> M4: Set {./relative-path} to external files for this module <<<

\def \modulefiles {./02-RISCVTRACENC}
% To add an external file, use:
% (Replace '---' with actual filename)
% '\input{\modulefiles/tables/---.tex}' for tables
% '\includegraphics{\modulefiles/figures/---}' for figures
% '\subfile{\modulefiles/---}' for register files

% >>> M1: This module is in global project? <<<
% 'YES' - uncomment; 'NO' - comment out
\globaltrue

\ifglobal
    \documentclass[main\_\_CN.tex]{subfiles}
   
\else
    % Font "FandolSong-Regular" used by ctex does not have GSUB table.
    % It causes warnings that do not affect compilation and can be ignored.
    \PassOptionsToPackage{quiet}{fontspec}

    \documentclass[openany, 10pt]{book}
   
    % >>> M2: In ./00-shared/preamble-trm-module, also find
    %         settings for visibility of todo notes, watermark <<< 
    \usepackage{preamble-trm-module}


    % Enable Chinese typesetting
    \usepackage{ctex}

    % >>> M3: Temporary labels in Chinese
    %         you can add your own labels there <<<
    \myexternaldocument{./00-shared/temp-labels-cn}

\fi %\ifglobal


\setcounter{secnumdepth}{4}


% >>> M5: If this module needs any updates in
% './00-shared/preamble-trm-module.sty'
% (include additional package, etc.), contact your module PM
% or create the file './00-shared/changelog.tex' make the
% necessary updates yourself and document them in this file <<<


\begin{document}


% Sets language into which pre-defined bits of text are translated
\selectlanguage{chinese}

% Do not delete this block
\ifglobal
% Add the following front matter if
% this module is outside of global project
\else
    \InputIfFileExists{readme}{}{}
    {\let\clearpage\relax \listoftodos}
    {\let\clearpage\relax \listofdonetodos}
    \tableofcontents
    %\listoftables
    %\listoffigures
\fi


%%%%%%%%%%%%%%%%%%%%%%%%%
%%% TRM Module Begins %%%
%%%%%%%%%%%%%%%%%%%%%%%%%

\hypertarget{riscvtracenc}{}
\chapter{RISC-V 追踪编码器 (TRACE)}
\label{mod:riscvtracenc}


\chipname{} 的高性能 CPU (HP CPU) 通过追踪编码器 (Trace Encoder) 实现了对指令追踪接口的支持。追踪编码器连接 HP CPU 的指令追踪接口，将信息压缩成更小的数据包，存入内部 SRAM（详见章节 \ref{mod:sysmem} \textit{\nameref{mod:sysmem}}）中。

\begin{figure}[H]
        \centering
        \includegraphics[width=0.8\textwidth]{\modulefiles/figures/trace.pdf}
        \caption{追踪编码器概述图}
        \label{fig:encoder}
\end{figure}

\section{术语}

为了更好地说明 RISC-V 追踪编码器的功能，本章使用了以下术语。 

\begin{longtable}[c]{ R{5cm} p{10cm} }
\textbf{hart} &  RISC-V 硬件线程 \\
\textbf{分支 (branch)} & 在一定条件下改变执行流程的指令 \\
\textbf{无法推测的不连续地址 (uninferable discontinuity)} & 无法仅通过程序二进制文件推测出的程序计数器地址偏移\\
\textbf{delta} & 指令在存储器中的存储地址不连续时，程序计数器的地址偏移\\
\textbf{陷阱 (trap)} & 即异常或中断，会触发陷阱处理程序 \\
\textbf{筛选 (qualification)} & 指令符合筛选条件，即称为通过筛选，才会被追踪 \\
\textbf{te\_inst} & 编码器发送的数据包的名称 \\
\textbf{退役 (retire)} & 机器状态更新时，执行指令的最后一个阶段 \\

\end{longtable}

%%%%%%%%%%%%%%%%%%%%
%%% Introduction %%%
%%%%%%%%%%%%%%%%%%%%

\section{介绍}

对于复杂的系统来说，了解程序执行流程并非易事。造成软件不按预期工作的可能有很多种因素，比如和其他核、外设、实时事件的交互，执行不畅，或以上所有因素的结合。

调试器很难实时监测系统运行中的程序执行流程，因为其是侵入性的，会破坏系统的运行状态。但是，监测程序执行又很重要。

此时，可使用指令追踪来追踪程序执行。


\begin{figure}[H]
        \centering
        \includegraphics[width=1\textwidth]{\modulefiles/figures/trace_over.pdf}
        \caption{指令追踪概览}
        \label{fig:trace}
\end{figure}
 
图 \ref{fig:trace}  展示了指令追踪的原理图：
\begin{itemize}
    \item HP CPU 核有一个指令追踪接口，可以输出 HP CPU 执行指令的信息，包括指令地址、指令类型等。更多关于 \chipname{} HP CPU 指令追踪接口的信息，可参考章节 \ref{mod:riscvcpu} \textit{\nameref{mod:riscvcpu}}。  
    \item 追踪编码器连接 HP CPU 的指令追踪接口，将信息压缩成带宽更小的数据包，存入系统存储器中。
    \item 调试器可以将系统存储器中的追踪数据包通过 JTAG 或 USB Serial/JTAG 导出，然后用解码器解压缩、重建程序的执行流。追踪解码器通常是外部 PC 上的软件，在获取追踪数据包后，根据在 hart 上运行的程序二进制数据，重建程序的指令流。解码可以离线进行，也可以在 hart 执行时实时进行。
\end{itemize}

本章节将主要介绍 \chipname{} 追踪编码器的具体实现。

%%%%%%%%%%%%%%%%%%%%
%%% Feature List %%%
%%%%%%%%%%%%%%%%%%%%

\section{特性}
\label{sec:riscvtracenc-features}

\begin{itemize}
    \item 兼容 RISC-V Processor Trace 1.0（RISC-V 追踪规范 v1.0），实现的参数详见表 \ref{tab:encoder_parameters}
    \item 支持任意地址范围用作追踪存储器
    \item 具有两个同步模式：
        \begin{itemize}
            \item 同步计数器按包计数
            \item 同步计数器按周期计数
        \end{itemize}
    \item 支持丢包状态标识
    \item 支持丢包后自动重启
    \item 写追踪存储器时支持循环和非循环模式
    \item 具有两个中断：
        \begin{itemize}
            \item 包的大小超过配置的存储器空间时触发中断
            \item 丢包时触发中断
        \end{itemize}
    \item 具有 128 × 8 位 FIFO，用于缓存数据包
\end{itemize}

\begin{longtable}[c]{ |l|r|l| }
    \caption{追踪编码器参数}
    \label{tab:encoder_parameters}
    \\\hline\rowcolor{lightgray}
    \textbf{参数名称} & \textbf{值} & \textbf{描述} \\\hline
    \endfirsthead
    \textbf{参数名称} & \textbf{值} & \textbf{描述} \\\hline
    \endhead
    \endfoot
    \endlastfoot
 
    arch\_p & 0 & 追踪器符合RISC-V 追踪规范初始版本  \\\hline
    bpred\_size\_p & 0 & 不支持分支预测模式 \\\hline
    cache\_size\_p & 0 & 不支持跳转目标缓存模式 \\\hline
    call\_counter\_size\_p & 0 & 不支持隐式返回模式 \\\hline
    ctype\_width\_p & 0 & 输出包中不包含上下文信息 \\\hline
    context\_width\_p & 0 & 输出包中不包含上下文信息 \\\hline
    ecause\_width\_p & 5 & 异常原因宽度 \\\hline
    ecause\_choice\_p & 0 & 不支持特权和异常原因多选 \\\hline
    f0s\_width\_p & 0 & 不支持格式 0 数据包 \\\hline
    filter\_context\_p & 0 & \multirow{4}{*}{不支持筛选功能} \\\cline{1-2}
    filter\_excint\_p & 0 &  \\\cline{1-2}
    filter\_privilege\_p & 0 &  \\\cline{1-2}
    filter\_tval\_p & 0 &  \\\hline
    iaddress\_lsb\_p & 1 & 支持压缩指令 \\\hline
    iaddress\_width\_p & 32 & 指令总线为 32 位 \\\hline
    iretire\_width\_p & 1 & iretire 总线宽度 \\\hline
    ilastsize\_width\_p & 0 & ilastsize 总线宽度 \\\hline
    itype\_width\_p & 3 & itype 总线宽度  \\\hline
    noncontext\_p & 1 & 在 te\_inst 数据包中去掉上下文信息 \\\hline
    privilege\_width\_p & 1 & 仅支持机器和用户模式 \\\hline
    retires\_p & 1 & 每个块退役指令的最大数量 \\\hline
    return\_stack\_size\_p & 0 & 不支持隐式返回模式 \\\hline
    sijump\_p & 0 & 不支持根据顺序推测跳转 \\\hline
    taken\_branches\_p & 1 & 每个时钟周期仅一条指令退役 \\\hline
    impdef\_width\_p & 0 & 未实现 \\\hline
    
 \end{longtable}

上述参数的具体描述，详见 RISC-V Processor Trace Version 1.0 > 章节 Parameters and Discovery。

%%%%%%%%%%%%%%%%%%%%%%%%%%%%%%
%%% Architectural Overview %%%
%%%%%%%%%%%%%%%%%%%%%%%%%%%%%%

\section{架构概览}
\label{sec:riscvtracenc-arch-overview}


如图 \ref{fig:encoder} 所示，追踪编码器包含一个编码器、一个 FIFO、一个寄存器配置模块和一个发送控制模块。

编码器通过指令追踪接口接收 HP CPU 的指令信息，将指令信息压缩成不同的数据包，写入内部 FIFO。

发送控制模块通过 AHB 总线将 FIFO 中的数据写入内部 SRAM。

FIFO 的深度为 128，宽度为 8 位。存储器写入带宽不足时，FIFO 可能会溢出，发生丢包。如果发生丢包，编码器会发送一个包通知丢包情况，之后停止工作，直至 FIFO 为空。

%%%%%%%%%%%%%%%%%%%%%%%%%%%%%%
%%% Functional 描述 %%%
%%%%%%%%%%%%%%%%%%%%%%%%%%%%%%

\section{功能描述}
\label{sec:riscvtracenc-func-desc}

\subsection{同步}

为保障追踪的有效性，追踪过程中必须定期同步。同步可以通过发送完整的指令地址实现。同步计数器的值达到 \hyperref[regdesc:TRACERESYNCPROLONGEDREG] {TRACE\_RESYNC\_PROLONGED\_REG} 寄存器的 \hyperref[fielddesc:TRACERESYNCPROLONGED]{TRACE\_RESYNC\_PROLONGED} 字段的值时，编码器会发送一个同步数据包（格式 3 子格式 0，见章节 \ref{sec:sec:riscvtracenc-format-3-sub-0}）。

追踪编码器支持两种同步模式，可通过 \hyperref[fielddesc:TRACERESYNCMODE]{TRACE\_RESYNC\_MODE} 配置：
\begin{itemize}
    \item 0：同步计数器按周期计数
    \item 1：同步计数器按数据包计数
\end{itemize}

您可通过增加 \hyperref[regdesc:TRACERESYNCPROLONGEDREG]{TRACE\_RESYNC\_PROLONGED\_REG} 的值，降低发送同步数据包的频率，从而降低同步数据包对带宽的占用。

\subsection{锚定}

由于数据包的长度不定，为在写入存储器时分清数据包之间的边界，\chipname{} 通过在数据包之间插入 0 ，从而能够发现数据包头：
\begin{itemize}
    \item 数据包的长度最大为 13 字节，因此数据包中不能出现连续 14 个值为 0 的字节。连续 14 个及以上值为 0 的字节后，第一个不为 0 的字节肯定是数据包的第一个字节。 
    \item 每发送 128 个数据包，编码器会向存储器写入 14 个值为 0 的字节作为锚定符。
\end{itemize}


\subsection{写存储器模式}

写追踪存储器时，追踪数据包的大小可能会超过追踪存储器的容量。此时可配置写存储器的模式，选择是否循环写入数据包：

\begin{itemize}
    \item 循环模式：追踪数据包的大小超过追踪存储器的容量（即\hyperref[regdesc:TRACEMEMCURRENTADDRREG]{TRACE\_MEM\_CURRENT\_ADDR\_REG} 的值达到 \hyperref[regdesc:TRACEMEMENDADDRREG]{TRACE\_MEM\_END\_ADDR\_REG} 的值）时，存储器循环写入数据，回到存储器的起始地址 \hyperref[regdesc:TRACEMEMSTARTADDRREG]{TRACE\_MEM\_START\_ADDR\_REG}，此前写入的数据会被新数据覆盖。
    \item 非循环模式：追踪数据包的大小超过追踪存储器的容量时，存储器不循环，停在 \hyperref[regdesc:TRACEMEMENDADDRREG]{TRACE\_MEM\_END\_ADDR\_REG}，此前写入的数据保留。
\end{itemize}

\subsection{自动重启}

因 FIFO 溢出而丢包时，编码器将停止工作，需要通过软件再次开启。如果 \hyperref[regdesc:TRACETRIGGERREG]{TRACE\_TRIGGER\_REG} 寄存器的 \hyperref[fielddesc:TRACERESTARTENA]{TRACE\_RESTART\_ENA} 位为 1，FIFO 清空后编码器会立即自动重启，无需软件重启。

如果开启了自动重启，即使编码器已经停止追踪，也可能会自动重启。因此，要关闭编码器，必须先清除 \hyperref[regdesc:TRACETRIGGERREG]{TRACE\_TRIGGER\_REG} 寄存器的 \hyperref[fielddesc:TRACERESTARTENA]{TRACE\_RESTART\_ENA} 位从而关闭自动重启。

%%%%%%%%%%%%%%%%%%%%%%%%%%%%%%
%%% Encoder Output Packets %%%
%%%%%%%%%%%%%%%%%%%%%%%%%%%%%%

\section{编码器输出数据包}
\label{sec:riscvtracenc-packetes}

本节主要介绍 \chipname{} 追踪编码器输出的数据包格式。\chipname{} 仅实现了必需的指令 delta 追踪，不支持下列可选功能：

\begin{itemize}
    \item delta 地址模式（支持运行时间配置模式）
    \item 上下文 (context) 信息和所有上下文相关的字段
    \item 可选的边带信号 (Sideband Signals)
    \item 调试模块触发输出
\end{itemize}

更多关于上述功能的细节，请参考 RISC-V Processor Trace 1.0。


\begin{figure}[H]
        \centering
        \includegraphics[width=0.7\textwidth]{\modulefiles/figures/package.pdf}
        \caption{追踪数据包格式}
        \label{fig:encoder-data-packet}
    \end{figure}

追踪数据包包括头部、索引和有效载荷，依次按照头部、索引和有效载荷的顺序以比特流形式传输。头部、索引和有效载荷中的字段按照下文表格的罗列顺序依次传输，多位字段先传输最低有效位。



\subsection{头部}

头部为 1 个字节，格式如表 \ref{tab:header} 所示。


\begin{longtable}[c]{ |l|r|l|l| }
    \caption{头部格式}
    \label{tab:header}
    \\\hline\rowcolor{lightgray}
    \textbf{字段} & \textbf{位} & \textbf{描述} & \textbf{值}\\\hline
    \endfirsthead
    \textbf{字段} & \textbf{位} & \textbf{描述} & \textbf{值}\\\hline
    \endhead
    \endfoot
    \endlastfoot
 
    length & 5 & 数据包的总长度 & 4\verb+~+13 \\\hline
    placeholder & 3 & 保留 & 0 \\\hline
 \end{longtable}

    
\subsection{索引}

索引为 2 个字节，格式如表 \ref{tab:index} 所示。

\begin{longtable}[c]{ |l|r|l|l| }
    \caption{索引格式}
    \label{tab:index}
    \\\hline\rowcolor{lightgray}
    \textbf{字段} & \textbf{位} & \textbf{描述} & \textbf{值}\\\hline
    \endfirsthead
    \textbf{字段} & \textbf{位} & \textbf{描述} & \textbf{值}\\\hline
    \endhead
    \endfoot
    \endlastfoot
 
    index & 16 & 每个数据包的索引 & 0\verb+~+65536 \\\hline
 \end{longtable}

\subsection{有效载荷}

有效载荷的宽度为 1 至 10 字节。

\subsubsection{格式 3}

格式 3 的数据包用于同步和发送支持信息。格式 3 有四种子格式，\chipname{} 仅支持 3 种。

\subparagraph{格式 3 子格式 0 - 同步}\label{sec:sec:riscvtracenc-format-3-sub-0}

该格式的数据包包含解码器完全识别指令所需的全部信息。开始追踪第一条指令时（除非指令恰好是异常处理程序中的第一条指令），或同步定时器达到阈值触发同步时，编码器会发送该格式的数据包。有效载荷的长度为 5 个字节。

\begin{longtable}[c]{ |l|r|L{10cm}| }
    \caption{格式 3 子格式 0}
    \label{tab:f3s0}
    \\\hline\rowcolor{lightgray}
    \textbf{字段名} & \textbf{位} & \textbf{描述} \\\hline
    \endfirsthead
    \rowcolor{lightgray}
    \textbf{字段名} & \textbf{位} & \textbf{描述} \\\hline
    \endhead
    \endfoot
    \endlastfoot
 
    format & 2 & 11 (sync)：同步 \\\hline
    subformat & 2 & 00 (start)：开始追踪，或开始同步 \\\hline
    branch & 1 & \makecell[l]{0：地址指向分支指令，且分支跳转\\1：指令非分支指令，或分支没有跳转} \\\hline
    privilege & 1 & 所追踪指令的特权等级 \\\hline
    address & 31 & 指令的完整地址。字段值需左移 1 位，以复原字节地址 \\\hline
    sign\_extend & 3 & 保留 \\\hline
 \end{longtable}

\subparagraph{格式 3 子格式 1 - 异常}
该格式的数据包同样包含解码器完全识别指令所需的全部信息。发生异常或触发中断时，编码器会发送该格式的数据包。该格式的数据包中包含了异常原因、异常指令的“陷阱值”，陷阱处理程序的地址和异常指令的地址。有效载荷的长度为 10 个字节。

\begin{longtable}[c]{ |l|r|L{10cm}| }
    \caption{格式 3 子格式 1}
    \label{tab:f3s1}
    \\\hline\rowcolor{lightgray}
    \textbf{字段名} & \textbf{位} & \textbf{描述} \\\hline
    \endfirsthead
    \textbf{字段名} & \textbf{位} & \textbf{描述} \\\hline
    \endhead
    \endfoot
    \endlastfoot
 
    format & 2 & 11 (sync)：同步 \\\hline
    subformat & 2 & 01 (exception)：异常原因和陷阱处理程序的地址 \\\hline
    branch & 1 & \makecell[l]{0：地址指向分支指令，且分支跳转\\1：指令非分支指令，或分支没有跳转} \\\hline
    privilege & 1 & 所追踪指令的特权等级 \\\hline
    ecause & 5 & 异常原因 \\\hline
    interrupt & 1 & 中断 \\\hline
    address & 31 & 指令的完整地址。地址必须向左偏移 1 位，以复原字节地址 \\\hline
    tvalepc & 32 & ecause 为 2、interrupt 为 0 时是异常地址，否则为陷阱值 \\\hline
    sign\_extend & 6 & 保留 \\\hline
 \end{longtable}


%%%%%%%%%%%%%%%%%%%%%%%%%%%%%%%%%%%%%%%%%%%%%%%%%%%%%%%%%%%%

\subparagraph{格式 3 子格式 3 - 支持}

该格式的数据包包含解码器需要的辅助信息，在追踪结束时发出。有效载荷的长度为 1 个字节。

\begin{longtable}[c]{ |l|r|L{10cm}| }
    \caption{格式 3 子格式 3}
    \label{tab:f3s3}
    \\\hline\rowcolor{lightgray}
    \textbf{字段名} & \textbf{位} & \textbf{描述} \\\hline
    \endfirsthead
    \rowcolor{lightgray}
    \textbf{字段名} & \textbf{位} & \textbf{描述} \\\hline
    \endhead
    \endfoot
    \endlastfoot
 
    format & 2 & 11 (sync)：同步 \\\hline
    subformat & 2 & 11 (support)：解码器需要的辅助信息 \\\hline
    enable & 1 & 表示编码器已开启 \\\hline
    \multirow{5}{*}{qual\_status} & \multirow{5}{*}{2} & 表示指令的筛选状态：
    \begin{itemize}
        \item 00 (no\_change)：指令筛选状态没有变化 
        \item 01 (ended\_rep)：指令筛选已结束，上一个指令为最后一个通过筛选的指令 
        \item 10 (trace lost)：丢失了一个或多个包 
        \item 11 (ended\_upd)：指令筛选已结束，即便上一个 te\_inst 
    \end{itemize}
    不是最后一个通过筛选的指令，也会因指令地址不连续、无法推测而发送
    \\\hline
    sign\_extend & 1 & 保留 \\\hline

 \end{longtable}




\subsubsection{格式 2}

该格式的数据包仅包含指令地址，在必须注明指令地址时使用，包中没有分支信息。有效载荷的长度为 5 个字节。

\begin{longtable}[c]{ |l|r|L{10cm}| }
    \caption{格式 2}
    \label{tab:f2}
    \\\hline\rowcolor{lightgray}
    \textbf{字段名} & \textbf{位} & \textbf{描述} \\\hline
    \endfirsthead
    \textbf{字段名} & \textbf{位} & \textbf{描述} \\\hline
    \endhead
    \endfoot
    \endlastfoot
 
    format & 2 & 10 (addr-only)：无分支信息 \\\hline
    address & 31 & 指令的完整地址 \\\hline
    notify & 1 & \chipname{} 不支持通知功能，因此该位永远和地址的最高有效位一致 \\\hline
    updiscon & 1 & 若该位的值和 notify 的值不同，则表示指令的地址指向一个无法推测的不连续的目标地址，且该地址为异常、特权变更或同步之前的最后一条指令 \\\hline
    sign\_extend & 5
    &  \\\hline
 \end{longtable}


\subsubsection{格式 1}

该格式的数据包包含分支信息，在必须注明分支信息（比如分支映射已满时）或在上一个包后出现一个以上分支指令、必须注明指令地址时使用。仅支持完整地址模式。

\subparagraph{格式 1，有地址时 branch\_map 字段}

有效载荷长度不定。

\begin{longtable}[c]{ |l|r|L{10cm}| }
    \caption{格式 1，有地址}
    \label{tab:f1a}
    \\\hline\rowcolor{lightgray}
    \textbf{字段名} & \multicolumn{1}{c|}{\textbf{位}} & \textbf{描述} \\\hline
    \endfirsthead
    \rowcolor{lightgray}
    \textbf{字段名} & \multicolumn{1}{c|}{\textbf{位}} & \textbf{描述} \\\hline
    \endhead
    \endfoot
    \endlastfoot
 
    format & 2 & 01：包含分支信息 \\\hline
    
    \multirow{8}{*}{branches} & \multirow{8}{*}{5} & 该字段的值决定了 branch\_map 的位宽，具体如下：
    \begin{itemize}
        \item 0：（该格式中此值无效）
        \item 1：1 位 
        \item 2-3：3 位 
        \item 4-7：7 位 
        \item 8-15：15 位 
        \item 16-31：31 位 
    \end{itemize}
    比如，若 branches = 12，则 branch\_map 长度为 15 位，低 12 位有效。 
    \\\hline
    \multirow{5}{*}{branch\_map} & \multirow{5}{*}{不定} & 表示分支是否跳转的位数组。位 0 对应最早执行的分支指令。每个位的含义如下：
    \begin{itemize}
        \item 0：分支已跳转 
        \item 1：分支未跳转 
    \end{itemize}
    该字段的位宽不定，取决于 branches 字段的值 \\\hline
    address & 31 & 指令的完整地址 \\\hline
    notify & 1 & \chipname{} 不支持通知功能，因此该位永远和地址的最高有效位一致\\\hline
    updiscon & 1 & 若该位的值和 notify 的值不同，则表示指令的地址指向一个无法推测的不连续的目标地址，且该地址为异常、特权变更或同步之前的最后一条指令 \\\hline
    
    \multirow{6}{*}{sign\_extend} & \multirow{6}{*}{不定} & 该字段的位宽不定，取决于 branches 字段的值，具体如下： 
    \begin{itemize}
        \item 1：7 位  
        \item 2-3：5 位  
        \item 4-7：1 位  
        \item 8-15：1 位  
        \item 16-32：31 位 
    \vspace{-1.5em}
    \end{itemize}
    \\\hline
 \end{longtable}
 
 \subparagraph{格式 1， 无地址时 branch\_map 字段}
有效载荷长度为 5 个字节。
\begin{longtable}[c]{ |l|r|L{10cm}| }
    \caption{格式 1，无地址}
    \label{tab:f1}
    \\\hline\rowcolor{lightgray}
    \textbf{字段名} & \textbf{位} & \textbf{描述} \\\hline
    \endfirsthead
    \textbf{字段名} & \textbf{位} & \textbf{描述} \\\hline
    \endhead
    \endfoot
    \endlastfoot
 
    format & 2 & 01：包含分支信息 \\\hline
    branches & 5 & 该字段的值决定了 branch\_map 的位宽，仅 0 有效，对应 branch\_map 的位宽为 31 位 \\\hline
    
    \multirow{4}{*}{branch\_map} & 31 & 表示分支是否跳转的位数组。位 0 对应最早执行的分支指令。每个位的含义如下：
    \begin{itemize}
        \item 0：分支已跳转 
        \item 1：分支未跳转 
    \vspace{-1.5em}
    \end{itemize}
    \\\hline

    sign\_extend & 2 & 保留 \\\hline
 \end{longtable}

%%%%%%%%%%%%%%%%%
%%% Interrupt %%%
%%%%%%%%%%%%%%%%%

\section{中断}

\begin{itemize}
    \item TRACE\_MEM\_FULL\_INTR 中断：数据包大小超过追踪存储器容量（即 \hyperref[regdesc:TRACEMEMCURRENTADDRREG]{TRACE\_MEM\_CURRENT\_ADDR\_REG} 的值达到 \hyperref[regdesc:TRACEMEMENDADDRREG]{TRACE\_MEM\_END\_ADDR\_REG} 的值）时触发。如有必要，可使能中断通知 HP CPU 处理，比如可再申请新的存储器空间。
    \item TRACE\_FIFO\_OVERFLOW\_INTR 中断：内部 FIFO 溢出时触发，表示一个或多个数据包丢失。
\end{itemize}

使能追踪编码器中断后，如果要让 HP CPU 响应中断，需要通过中断矩阵将追踪编码器中断映射到 CPU 中断号，详情请参考章节 \ref{mod:intmtrx} \textit{\nameref{mod:intmtrx}}。

%%%%%%%%%%%%%%%%%%%%%%%%%%%%%%
%%% Programming Procedures %%%
%%%%%%%%%%%%%%%%%%%%%%%%%%%%%%

\section{编程流程}
\label{sec:riscvtracenc-prog-procedures}

\subsection{使能编码器}


\begin{itemize}
    \item 通过 \hyperref[regdesc:TRACEMEMSTARTADDRREG]{TRACE\_MEM\_START\_ADDR\_REG} 和 \hyperref[regdesc:TRACEMEMENDADDRREG]{TRACE\_MEM\_END\_ADDR\_REG} 配置追踪存储器的地址空间
    
    \item 向 \hyperref[fielddesc:TRACEMEMCURRENTADDRUPDATE]{TRACE\_MEM\_CURRENT\_ADDR\_UPDATE} 写 1，将 \hyperref[regdesc:TRACEMEMCURRENTADDRREG]{TRACE\_MEM\_CURRENT\_ADDR\_REG} 的值更新为 \hyperref[regdesc:TRACEMEMSTARTADDRREG]{TRACE\_MEM\_START\_ADDR\_REG} 的值
    
    \item（可选）通过 \hyperref[regdesc:TRACETRIGGERREG]{TRACE\_TRIGGER\_REG} 寄存器的 \hyperref[fielddesc:TRACEMEMLOOP]{TRACE\_MEM\_LOOP} 位配置写存储器的模式
        \begin{itemize}
            \item 0：非循环模式
            \item 1：循环模式（默认）
        \end{itemize}
        
    \item 通过 \hyperref[regdesc:TRACERESYNCPROLONGEDREG]{TRACE\_RESYNC\_PROLONGED\_REG} 寄存器的 \hyperref[fielddesc:TRACERESYNCMODE]{TRACE\_RESYNC\_MODE} 位配置同步模式
        \begin{itemize}
                    \item 0：同步计数器按周期计数（默认）
                    \item 1：同步计数器按数据包计数
                \end{itemize}
    
    \item（可选）通过 \hyperref[regdesc:TRACERESYNCPROLONGEDREG]{TRACE\_RESYNC\_PROLONGED\_REG} 配置同步计数器的阈值，默认为 128
    
    \item（可选）使能中断
        \begin{itemize}
            \item 向 \hyperref[regdesc:TRACEINTRENAREG]{TRACE\_INTR\_ENA\_REG} 寄存器的相应位写 1 开启对应中断
            \item 向\hyperref[regdesc:TRACEINTRCLRREG]{TRACE\_INTR\_CLR\_REG} 寄存器的相应位写 1 清除对应中断
            \item 读 \hyperref[regdesc:TRACEINTRRAWREG]{TRACE\_INTR\_RAW\_REG} 寄存器的值，获取中断状态
        \end{itemize}
        
    \item （可选）向 \hyperref[regdesc:TRACETRIGGERREG]{TRACE\_TRIGGER\_REG} 寄存器的 \hyperref[fielddesc:TRACERESTARTENA]{TRACE\_RESTART\_ENA} 位写 1，使能自动重启。该功能默认开启
    
    \item 向 \hyperref[regdesc:TRACETRIGGERREG]{TRACE\_TRIGGER\_REG} 寄存器的 \hyperref[fielddesc:TRACETRIGGERON]{TRACE\_TRIGGER\_ON} 字段写 1，使能编码器
    
\end{itemize}

编码器使能后，会一直追踪 HP CPU 的指令接口，向追踪存储器写入数据包。

\subsection{关闭编码器}


\begin{itemize}
    \item 通过清除 \hyperref[regdesc:TRACETRIGGERREG]{TRACE\_TRIGGER\_REG} 寄存器的 \hyperref[fielddesc:TRACERESTARTENA]{TRACE\_RESTART\_ENA} 位关闭自动重启
    
    \item 向 \hyperref[regdesc:TRACETRIGGERREG]{TRACE\_TRIGGER\_REG} 寄存器的 \hyperref[fielddesc:TRACETRIGGEROFF]{TRACE\_TRIGGER\_OFF} 位写 1 关闭编码器
    
    \item 读取 \hyperref[fielddesc:TRACEFIFOEMPTY]{TRACE\_FIFO\_EMPTY} 位的值，确认 FIFO 中数据是否都已写入追踪存储器

\end{itemize}

\subsection{解码数据包}

\begin{itemize}
    \item 找到解码的第一个地址
        \begin{itemize}
            \item 读 \hyperref[regdesc:TRACEINTRRAWREG]{TRACE\_INTR\_RAW\_REG} 寄存器的 \hyperref[fielddesc:TRACEMEMFULLINTRRAW]{TRACE\_MEM\_FULL\_INTR\_RAW} 位，查看追踪存储器是否已满
            \begin{itemize}
                \item 如果值为 0，从 \hyperref[regdesc:TRACEMEMSTARTADDRREG]{TRACE\_MEM\_START\_ADDR\_REG} 寄存器读取追踪数据包
                \item 如果值为 1，且开启了存储器循环模式，则之前的数据包已被覆盖。此时，读取 \hyperref[regdesc:TRACEMEMCURRENTADDRREG]{TRACE\_MEM\_CURRENT\_ADDR\_REG} 的值，获取最后一个写入的地址，该地址即为解码器处理的第一个地址
            \end{itemize}
        \end{itemize}
    
    \item 用解码器解码数据包
        \begin{itemize}
            \item 解码器从第一个地址开始读取所有数据包，结合二进制文件重建数组
            \item 如章节 \ref{sec:riscvtracenc-packetes} 所述，编码器向存储器分区边界写 14 个值为 0 的字节。因此出现 14 个值为 0 的字节时，则表示下一个非 0 字节为新数据包的头部
        \end{itemize}
\end{itemize}




% >>> If your module has no registers, delete sample text
%       for Sections Register Summary and Registers below <<<

%%%%%%%%%%%%%%%%%%%%%%%%
%%% Register Summary %%%
%%%%%%%%%%%%%%%%%%%%%%%%

% >>> Update the sample text below and insert
%     the info on register summary into the subfile <<<

\clearpage
\hypertarget{riscvtracenc-reg-summ}{}
\section{寄存器列表}
\label{sec:riscvtracenc-reg-summ}

本小节的所有地址均为相对于 RISC-V 追踪编码器基地址的地址偏移量（相对地址），具体基地址请见章节 \ref{mod:sysmem} \textit{\nameref{mod:sysmem}} 中的表 \ref{tab:sysmem-base-address}。

请查看章节 \hyperref[glossary-access-types]{\textit{寄存器的访问类型}}，了解“访问”列缩写的含义。

\subfile{\modulefiles/02-RISCVTRACENC-reg-summ__CN}


%%%%%%%%%%%%%%%%%
%%% Registers %%%
%%%%%%%%%%%%%%%%%

% >>> Update the sample text below and insert
%      the info on registers into the subfile <<<

\clearpage
\section{寄存器}

本小节的所有地址均为相对于 RISC-V 追踪编码器基地址的地址偏移量（相对地址），具体基地址请见章节 \ref{mod:sysmem} \textit{\nameref{mod:sysmem}} 中的表 \ref{tab:sysmem-base-address}。

\subfile{\modulefiles/02-RISCVTRACENC-reg__CN}


\end{document}
